% =====================================================================
%  1 Technologieübersicht  (Aufgabenstellung: „Beschreibung und Begründung
%  der Wahl der verwendeten Technologien und Werkzeuge (z. B.
%  Programmiersprache, Frameworks, Bibliotheken)“)
%  Pflicht: Jede Third-Party-Bibliothek dokumentieren UND begründen,
%  inkl. Version und Lizenz. Alles, was über Python-Interpreter und
%  Standardbibliotheken hinausgeht, muss in der Auslieferung enthalten
%  sein (Aufgabenstellung 1) -> in der Benutzeranleitung beschreiben.
%  Wenig bekannte Software/Bibliotheken ins Literaturverzeichnis
%  (Zitierleitfaden Kap. 2.3.9).
% =====================================================================
\section{Technologieübersicht}
\TODO{Kurze Hinführung}

\subsection{Programmiersprache und Laufzeitumgebung}
\TODO{Python: Version, Begründung (objektorientiert, Vorgabe der Aufgabenstellung erfüllt), Windows 10/11}

\subsection{Bibliotheken und Frameworks}
\begin{longtable}{@{}L{3.0cm}L{1.6cm}L{2.2cm}L{3.4cm}L{5.0cm}@{}}
  \caption{Verwendete Bibliotheken und Frameworks}\label{tab:bibliotheken}\\
  \toprule
  \textbf{Name} & \textbf{Version} & \textbf{Lizenz} & \textbf{Zweck} & \textbf{Begründung} \\
  \midrule
  \endfirsthead
  \toprule
  \textbf{Name} & \textbf{Version} & \textbf{Lizenz} & \textbf{Zweck} & \textbf{Begründung} \\
  \midrule
  \endhead
  \bottomrule
  \endfoot
  \TODO{GUI-Framework} & \TODO{} & \TODO{} & Grafische Oberfläche & \TODO{Alternativen abgewogen} \\
  \TODO{} & \TODO{} & \TODO{} & \TODO{} & \TODO{} \\
\end{longtable}
\quelle{Eigene Darstellung.}

\subsection{Entwicklungswerkzeuge}
\TODO{IDE, Git/GitHub, Test-Framework, Paketierung (z. B. zur .exe), LaTeX, draw.io – jeweils kurz begründen}
